%
% 12-schwingungen.tex -- Schwingungen
%
% (c) 2026 Prof Dr Andreas Müller
%
\subsection{Schwingungen}
Die Schwingungsdifferentialgleichung ist eine Differentialgleichung
der Form
\begin{equation}
y''
=
-\omega^2\,y.
\label{buch:differentialgleichungen:potenzreihen:eqn:schwingungsdgl}
\end{equation}
Auch in diesem Fall sind die Koeffizienten konstant und somit
analytisch und es ist gerechtfertigt, eine analytische Lösung
zu erwarten und mit einem Potenzreihenansatz zu suchen.

%
% Reell analytische Lösung
%
\subsubsection{Reell analytische Lösung}
Die zweite Ableitung des
Potenzreihenansatzes~\eqref{buch:differentialgleichungen:potenzreihen:eqn:ansatz}
ist
\[
y''(x)
=
\sum_{l=2}^\infty
a_l l(l-1) x^{l-2}
.
\]
Setzen wir dies in die Differentialgleichung
\eqref{buch:differentialgleichungen:potenzreihen:eqn:schwingungsdgl}
ein, ergibt sich
\begin{align*}
\sum_{l=2}^\infty
a_l l(l-1) x^{l-2}
&=
-\omega^2
\sum_{l=0}^\infty
a_lx^l
\intertext{Auch in diesem Fall ist es notwendig, auf der rechten
Seite $l+2$ durch $l$ zu ersetzen, was auf die Gleichung}
\sum_{l=0}^\infty
a_{l+2} (l+1)(l+2) x^{l}
&=
-\omega^2
\sum_{l=0}^\infty
a_lx^l
\end{align*}
Koeffizientenvergleich ergibt die Rekursionsformel
\begin{align*}
(l+1)(l+2)
a_{l+2} 
&=
-\omega^2
a_l
\intertext{für $l>0$ oder}
a_{l+2}
&=
-\frac{\omega^2}{(l+1)(l+2)} a_{l}.
\end{align*}
Man beachte, dass die Rekursionsformel nur entweder gerade oder ungerade
indizierte Koeffizienten untereinander in Beziehung setzt.
Die gerade indizierten Koeffizienten müssen sich daher alle durch
$a_0$ ausdrücken lassen, während die ungerade indizierten Koeffizienten
von $a_1$ abhängen.
Die Lösungsfunktion ist also erst durch die Angabe von $a_0=y(0)$ und
$a_1=y'(0)$ eindeutig festgelegt.

%
% Lösung der Rekursionsgleichung
%
\subsubsection{Lösung der Rekursionsgleichung}
Diese Rekursionsformeln ermöglichen, die Koeffizienten $a_l$ mit $l\ge 2$
durch $a_0$ und $a_1$ als
\begin{align*}
a_2
&=
-\omega^2 \frac{1}{1\cdot 2} a_0
\\
a_3
&=
-\omega^2 \frac{1}{2\cdot 3} a_1
%=
%-\omega^2 \frac{1}{3!} a_1
\\
a_4
&=
-\omega^2 \frac{1}{3\cdot 4} a_2
=
\phantom{-}\omega^4 \frac{1}{4!} a_0
\\
a_5
&=
-\omega^2 \frac{1}{4\cdot 5} a_3
=
\phantom{-}\omega^4 \frac{1}{5!} a_1
\\
a_6
&=
-\omega^2 \frac{1}{5\cdot 6} a_4
=
-\omega^6 \frac{1}{6!} a_0
\\
a_7
&=
-\omega^2 \frac{1}{6\cdot 7} a_5
=
-\omega^6 \frac{1}{7!} a_1
\intertext{auszudrücken oder allgemein}
a_{2n}   &= (-1)^n \omega^{2n} \frac{1}{(2n)!}   a_0
\\
a_{2n+1} &= (-1)^n \omega^{2n} \frac{1}{(2n+1)!} a_1.
\end{align*}
Daraus kann die Lösung als
\begin{align}
y(x)
&=
a_0
\sum_{n=0}^\infty
(-1)^n\omega^{2n}\frac{1}{(2n)!} x^{2n}
+
a_1
\sum_{n=0}^\infty
(-1)^n\omega^{2n}\frac{1}{(2n+1)!} x^{2n+1}
\notag
\\
&=
a_0
\sum_{n=0}^\infty
\frac{(-1)^n}{(2n)!} (\omega x)^{2n}
+
\frac{a_1}{\omega}
\sum_{n=0}^\infty
\frac{(-1)^n}{(2n+1)!} (\omega x)^{2n+1}
\notag
\\
&=
a_0\cos(\omega x)
+
\frac{a_1}{\omega}\sin(\omega x)
\label{buch:differentialgleichungen:potenzreihen:eqn:schwloesung}
\end{align}
ableiten.
Die Lösungen sind also Linearkombinationen von Sinus-
bzw.~Kosinus-Schwin\-gun\-gen.

Auch in diesem Fall ist die zweite Ableitung der Lösung
\eqref{buch:differentialgleichungen:potenzreihen:eqn:schwloesung}
\[
y''(x)
=
-\omega^2 a_0\cos(\omega x)
-
\omega^2 \frac{a_1}{\omega} \sin(\omega x)
=
-\omega^2 y(x),
\]
und erfüllt damit wie erwartet die Schwingungsdifferentialgleichung.

%
% Komplexe Lösungen
%
\subsubsection{Komplexe Lösungen}
Die Differentialgleichung 
\eqref{buch:differentialgleichungen:potenzreihen:eqn:schwingungsdgl}
kann man auch in der Form
\begin{equation}
\biggl(\frac{d}{dx}-i\omega\biggr)
\biggl(\frac{d}{dx}+i\omega\biggr)
y
=
0
\label{buch:differentialgleichungen:potenzreihen:eqn:faktschwing}
\end{equation}
schreiben.
Lösungen der Differentialgleichungen erster Ordnung
\begin{equation}
\biggl(\frac{d}{dx} - i \omega\biggr)y = 0
\qquad\text{und}\qquad
\biggl(\frac{d}{dx} + i \omega\biggr)y = 0
\label{buch:differentialgleichungen:potenzreihen:eqn:einzelschwing}
\end{equation}
ergeben linear unabhängige Lösungen der Differentialgleichung
\eqref{buch:differentialgleichungen:potenzreihen:eqn:faktschwing}.
Die Differentialgleichungen
\eqref{buch:differentialgleichungen:potenzreihen:eqn:einzelschwing}
sind komplexe Differentialgleichungen der Form
\eqref{buch:differentialgleichungen:potenzreihen:eqn:expdgl}
und können daher mit der Exponentialreihe gelöst werden.
Die Lösungen von 
\eqref{buch:differentialgleichungen:potenzreihen:eqn:einzelschwing}
\begin{equation}
y_+(x) = e^{i\omega x}
\qquad\text{und}\qquad
y_-(x) = e^{-i\omega x}
\end{equation}
können zu den allgemeinen Lösungen der Schwingungsdifferentialgleichungen
linear kombiniert werden.
Die trigonometrischen Funktionen als Lösungen entstehen als die
Linearkombinationen
\[
\cos \omega x
=
\frac{y_+(x)+y_-(x)}{2}
\qquad\text{und}\qquad
\sin \omega x
=
\frac{y_+(x)-y_-(x)}{2i}.
\]
